% chap_01.tex
\chapter{General Project Framework}

\section*{Introduction}
\addcontentsline{toc}{section}{Introduction}

This chapter presents the general framework of the end-of-study project. We begin by describing the host organization, its environment, and its fields of activity. We then present the project context along with a critical study of existing human resources management solutions. Finally, we formulate the identified problem and present the proposed solution.

% ==================================================
\section{Project Context}

As part of obtaining the engineering degree in Computer Science, specializing in Web Technologies and Numerical Intelligence (TWIN), at the Private Higher School of Engineering and Technology (ESPRIT), during the 2025/2026 academic year, we carried out our end-of-study project within the company \textbf{TAC-TIC}.

This project is part of a broader effort to improve and digitalize internal processes, particularly those related to human resources management. Its aim is to design and develop an intelligent system capable of centralizing, automating, and optimizing HR activities within the company.

% ==================================================
\section{Presentation of the Host Company}

\subsection{General Overview}

Founded in 2012, \textbf{TAC-TIC} is a company specialized in IT services and innovative technological solutions. Since its creation, the company has aimed to position itself as a leading technology partner for its clients by offering a wide range of services covering IT infrastructure, cloud services, software development, and consulting.

TAC-TIC stands out for its commitment to quality, innovation, and customer satisfaction. It aims to become one of the most reputable IT service providers internationally by relying on the expertise of its teams and the trust of its partners.

\subsection{Fields of Activity}

The main activities of TAC-TIC include:

- IT infrastructure services,\\
- Cloud and virtualization solutions,\\
- Web and mobile application development,\\
- Consulting and support for digital transformation,\\
- DevOps practices and automation,\\
- Solutions integrating artificial intelligence and data analysis.

\subsection{Internal Organization}

The company adopts a structured organization built around several departments: general management, technical teams, project management, and support functions (human resources, finance, and administration). This organization enables effective project coordination and optimized resource management.

% ==================================================
\section{Project Presentation}

The project consists of designing an \textbf{Intelligent Human Resources Management System} intended to digitalize and optimize the company’s HR processes.

The main objective is to centralize the management of employees, attendance tracking, leave, and payroll within a single platform, while integrating an artificial intelligence module designed to generate performance indicators, predict absences, and support decision-making.

This solution is part of a digital transformation approach aimed at improving data reliability, action traceability, and operational efficiency.

% ==================================================
\section{Study of Existing Solutions}

\subsection{Current Context}

In many companies, human resources management relies on disparate tools or partially manual processes. The systems used are often limited to basic administrative functionalities, without real integration or advanced analytical capabilities.

\subsection{Critical Analysis}

The main limitations of traditional solutions are as follows:

- \textbf{Data fragmentation:} use of multiple non-integrated tools.\\
- \textbf{Manual processing:} payroll calculations and absence tracking require repetitive human intervention.\\
- \textbf{Lack of traceability:} difficulty in tracking the history of modifications and decisions.\\
- \textbf{Lack of analytical indicators:} few or no decision-making dashboards.\\
- \textbf{No predictive capability:} lack of advanced analysis mechanisms to anticipate absences or optimize resource allocation.

These limitations reduce the company’s ability to efficiently exploit its HR data and make strategic decisions based on objective analysis.

% ==================================================
\section{Comparative Table of Existing Tools}

The following table presents some existing HR solutions as well as their main limitations compared with our proposed solution.

\begin{table}[H]
\centering
\renewcommand{\arraystretch}{1.5}
\begin{tabular}{|p{3cm}|p{3cm}|>{\raggedright\arraybackslash}p{7cm}|}
\hline
\textbf{Tools} & \textbf{Logo} & \textbf{Limitations} \\ 
\hline

Odoo HR
& \makebox[3cm][c]{\includegraphics[width=2.2cm]{tpl/img/logo-odoo.png}}
& Complex configuration, advanced paid modules, no native predictive AI.
\\ \hline

Zoho People
& \makebox[3cm][c]{\includegraphics[width=2.2cm]{tpl/img/zoho people.png}}
& Limited payroll module depending on the region, mainly descriptive reporting, lack of intelligent prediction.
\\ \hline

BambooHR
& \makebox[3cm][c]{\includegraphics[width=2.2cm]{tpl/img/bamboo.png}}
& Payroll not available in all countries, limited customization, basic analytical features.
\\ \hline

\end{tabular}
\caption{Comparison of existing HR solutions}
\label{tab:comparatif-outils}
\end{table}

% ==================================================
\section{Problem Statement}

Given the increasing complexity of HR processes and the need to improve organizational performance, it has become essential to implement an integrated and intelligent solution.

The main problem can be formulated as follows:

\textit{How can we design and implement an intelligent human resources management system capable of automating administrative processes, ensuring data reliability, and introducing decision support based on analysis and prediction?}

% ==================================================
\section{Proposed Solution}

To address this problem, we propose the development of an intelligent system integrating:

- Centralized management of employees and their professional information,\\
- An attendance management module with presence and absence tracking,\\
- A leave management system integrated with the other modules,\\
- An automatic payroll calculation engine based on attendance data,\\
- An artificial intelligence module enabling KPI generation, absence prediction, and profile matching.

This solution aims to improve transparency, performance, and the quality of decision-making within the company.

% ==================================================
\section*{Conclusion}

In this chapter, we presented the general framework of the project as well as the host company. We analyzed the current state of human resources management systems and identified their main limitations. This analysis made it possible to formulate the project problem and justify the proposed solution. The next chapter will be devoted to the project initiation phase (Sprint 0), including the detailed analysis of requirements and the overall system design.